\chapter{Interpolimi, regresi dhe përshtatja e modeleve}
\begin{qellime}
Studenti ndërton polinomin e Lagranzhit dhe interpolimin Hermitian; dallon interpolimin nga regresi; zbaton katrorët më të vegjël linearë e jolinearë; interpreton mbetjet, peshat dhe pacaktueshmërinë e parametrave.
\end{qellime}

\section{Interpolimi polinomial}
Për nyje të dallueshme $(x_i,y_i)$ ekziston një polinom unik me gradë jo më të madhe se $n$ që kalon në $n+1$ pika. Forma e Lagranzhit është
\begin{equation}
 p_n(x)=\sum_{i=0}^{n}y_iL_i(x),\qquad
 L_i(x)=\prod_{j\ne i}\frac{x-x_j}{x_i-x_j}.
\end{equation}
Gabimi për $f\in C^{n+1}$ është
\begin{equation}
 f(x)-p_n(x)=\frac{f^{(n+1)}(\xi)}{(n+1)!}\prod_{i=0}^{n}(x-x_i).
\end{equation}
Formula tregon rolin e lëmueshmërisë dhe të vendosjes së nyjeve. Rritja e gradës në nyje të baraslarguara mund të përkeqësojë përafrimin pranë skajeve: fenomeni Runge.

\begin{figure}[ht]
\centering\includegraphics[width=.95\textwidth]{06_interpolimi_runge.pdf}
\caption{Interpolimi i funksionit Runge. Nyjet e Çebishevit zvogëlojnë lëkundjet në skaje.}
\end{figure}

Në llogaritje përdoret forma baricentrike, jo prodhimi i drejtpërdrejtë i polinomeve bazë. Për grupe të mëdha të dhënash preferohen splinet kubike, të cilat përdorin polinome lokale dhe ruajnë vazhdimësinë e dy derivateve të para.

\section{Interpolimi Hermitian dhe ekstrapolimi}
Interpolimi Hermitian përdor edhe vlerat e derivateve. Për dy pika, polinomi kubik përcaktohet nga $f(x_0)$, $f'(x_0)$, $f(x_1)$ dhe $f'(x_1)$. Kjo është bazë për trajektore të lëmuara dhe skema integrimi.

Ekstrapolimi vlerëson jashtë intervalit të të dhënave dhe është shumë më i rrezikshëm. Ekstrapolimi Richardson shfrytëzon një zhvillim të njohur të gabimit,
\begin{equation}
 A(h)=A+c h^p+\Oh(h^{p+1}),qquad
 A\approx\frac{2^pA(h/2)-A(h)}{2^p-1}.
\end{equation}
Ai përdoret për të rritur rendin ose për të vlerësuar kufirin $h\to0$.

\section{Regresi linear}
Në regres, pikat nuk kërkohet të kalohen saktësisht. Modeli $\vect y=X\vect\beta+\vect\varepsilon$ përshtatet duke minimizuar
\begin{equation}
 S(\vect\beta)=\|X\vect\beta-\vect y\|_2^2.
\end{equation}
Ekuacionet normale janë $X^TX\hat{\vect\beta}=X^T\vect y$, por formimi i $X^TX$ katrorëzon numrin e kushtëzimit. Faktorizimi QR është më i qëndrueshëm. Kur devijimet standarde $\sigma_i$ njihen, minimizohet
\begin{equation}
 \chi^2=\sum_i\frac{[y_i-f(x_i;\vect\beta)]^2}{\sigma_i^2}.
\end{equation}

\begin{figure}[ht]
\centering\includegraphics[width=.72\textwidth]{07_regresioni_kepler.pdf}
\caption{Linearizimi logaritmik i $T^2\propto a^3$. Pjerrësia e pritur për $\ln T$ kundrejt $\ln a$ është $3/2$.}
\end{figure}

\begin{shembull}[Ligji i tretë i Keplerit]
Modeli $T=C a^p$ bëhet $\ln T=\ln C+p\ln a$. Përshtatja jep $p$ dhe pacaktueshmërinë e tij. Transformimi logaritmik ndryshon strukturën e gabimit; nëse gabimi është aditiv në $T$, përshtatja jolineare në shkallën origjinale mund të jetë më e përshtatshme.
\end{shembull}

\section{Mbetjet dhe diagnostika}
Mbetjet $r_i=y_i-\hat y_i$ nuk janë mbeturina; ato testojnë modelin. Një strukturë sistematike në $r_i$ tregon jolinearitet të munguar, heteroskedasticitet, korrelacion ose pikë jashtë modelit. Koeficienti $R^2$ nuk mjafton për vleftësim dhe nuk duhet përdorur si provë shkakësie.

Për gabime Gauss të pavarura me variancë $\sigma^2$, kovarianca e parametrave është afërsisht
\begin{equation}
 \Cov(\hat{\vect\beta})=\sigma^2(X^TX)^{-1}.
\end{equation}
Në modele jolineare përdoret Jakobiani në optimum, bootstrap-i ose kampionimi probabilitar.

\section{Shmangiet minimale absolute}
Minimizimi $\sum_i|r_i|$ është më robust ndaj pikave ekstreme se katrorët më të vegjël. Funksioni nuk është i derivueshëm në zero, por problemi mund të formulohet si program linear. Zgjedhja e humbjes shpreh një supozim statistik: $L_2$ lidhet me gabimin Gauss, ndërsa $L_1$ me shpërndarje me bishta më të rëndë.

\begin{lidhje}
Laboratori përdor të dhëna planetare, përshtat $T=C a^p$, vizaton mbetjet dhe krahason regresin e papeshuar, të peshuar dhe një përshtatje robuste.
\end{lidhje}

\section*{Pyetje kontrolli}
\begin{enumerate}
\item Pse interpolimi me gradë të lartë mund të dështojë edhe për funksion të lëmuar?
\item Kur linearizimi logaritmik ndryshon interpretimin statistik të gabimit?
\item Çfarë informacioni japin mbetjet që nuk e jep $R^2$?
\end{enumerate}
\section{Algoritmet e interpolimit dhe përshtatjes}
\subsection{Interpolimi baricentrik}
Forma baricentrike është më e qëndrueshme se vlerësimi i drejtpërdrejtë i bazës së Lagranzhit. Peshat janë
\begin{equation}
 w_i=\left[\prod_{j\ne i}(x_i-x_j)\right]^{-1},\qquad
 p(x)=\frac{\sum_iw_iy_i/(x-x_i)}{\sum_iw_i/(x-x_i)}.
\end{equation}
Nëse $x$ përputhet me një nyje, kthehet drejtpërdrejt $y_i$.

\begin{lstlisting}[language=Python,caption={Interpolimi baricentrik nga SciPy.}]
import numpy as np
from scipy.interpolate import BarycentricInterpolator

f = lambda x: 1.0/(1.0 + 25.0*x*x)
n = 15
i = np.arange(n)
x_cheb = np.cos((2*i + 1)*np.pi/(2*n))
interp = BarycentricInterpolator(x_cheb, f(x_cheb))

x = np.linspace(-1, 1, 1000)
gabimi = np.max(np.abs(interp(x) - f(x)))
print("gabimi maksimal =", gabimi)
\end{lstlisting}

\subsection{Katrorët më të vegjël me QR}
\begin{lstlisting}[language=Python,caption={Përshtatja e një modeli linear në parametra.}]
def pershtat_linearisht(X, y, sigma=None):
    if sigma is not None:
        W = 1.0 / np.asarray(sigma)
        X = X * W[:, None]
        y = y * W
    beta, mbetje, rank, singular = np.linalg.lstsq(X, y, rcond=None)
    return beta, mbetje, rank, singular

# ln(T) = beta0 + beta1 ln(a)
X = np.column_stack([np.ones_like(a), np.log(a)])
beta, *_ = pershtat_linearisht(X, np.log(T))
print("eksponenti i Keplerit =", beta[1])
\end{lstlisting}
\texttt{lstsq} përdor faktorizime numerikisht më të sigurta se zgjidhja eksplicite e ekuacioneve normale.

\subsection{Përshtatja jolineare}
\begin{lstlisting}[language=Python,caption={Përshtatja e modelit $T=C a^p$ në shkallën origjinale.}]
from scipy.optimize import curve_fit

def kepler(a, C, p):
    return C * a**p

param, kov = curve_fit(kepler, a, T, p0=(1.0, 1.5), sigma=sigma_T,
                       absolute_sigma=True)
gabim_param = np.sqrt(np.diag(kov))
print(param, gabim_param)
\end{lstlisting}

\subsection{Diagnostika e mbetjeve}
Pas përshtatjes vizatohen $r_i$ kundrejt $x_i$, histogrami dhe, kur rendi kohor ka rëndësi, autokorrelacioni. Një lakim sistematik tregon model të mangët; formë hinke tregon variancë jo të njëtrajtshme; pika me ndikim të madh kërkojnë kontroll fizik, jo fshirje automatike.

\begin{kujdes}
Termi \emph{përshtatje} përdoret për fitting, \emph{mbetje} për residual dhe \emph{shmangie minimale absolute} për least absolute deviations. \emph{Regres} i referohet modelit të statistikës, ndërsa \emph{interpolim} kalon saktësisht nëpër nyje.
\end{kujdes}

Për interpolimin, përafrimin dhe shembujt MATLAB/Python janë konsultuar \cite{hanelli,hoxha,datja}. Për terminologjinë e rizgjedhjes, mbetjeve dhe vlerësimit statistik janë përdorur edhe punimet doktorale të UT-së \cite{butka,tasho}.

\subsection{Interpolimi i Lagranzhit dhe termi i mbetjes}
Për nyje të ndryshme $x_0,\ldots,x_n$, polinomi interpolues është
\begin{equation}
 p_n(x)=\sum_{i=0}^{n}f(x_i)L_i(x),\qquad
 L_i(x)=\prod_{j\ne i}\frac{x-x_j}{x_i-x_j}.
\end{equation}
Vërtetimi është i drejtpërdrejtë: $L_i(x_j)=\delta_{ij}$, prandaj $p_n(x_j)=f(x_j)$. Uniciteti rrjedh sepse diferenca e dy polinomeve interpoluese do të kishte $n+1$ rrënjë, por gradë jo më të madhe se $n$.

Nëse $f\in C^{n+1}$, gabimi në një pikë $x$ është
\begin{equation}
 f(x)-p_n(x)=\frac{f^{(n+1)}(\xi_x)}{(n+1)!}
 \omega_{n+1}(x),\qquad
 \omega_{n+1}(x)=\prod_{i=0}^{n}(x-x_i).
\end{equation}
Kjo formulë ndan dy burime: lëmueshmërinë e funksionit dhe gjeometrinë e nyjeve. Nyjet e baraslarguara mund ta bëjnë $\omega_{n+1}$ shumë të madhe pranë skajeve; nyjet e Çebishevit e kontrollojnë këtë faktor dhe zbusin dukurinë Runge.

Forma baricentrike
\begin{equation}
 p_n(x)=\frac{\sum_i w_i f_i/(x-x_i)}{\sum_i w_i/(x-x_i)},
 \qquad w_i=\left[\prod_{j\ne i}(x_i-x_j)\right]^{-1},
\end{equation}
është numerikisht më e përshtatshme se zgjerimi i polinomit në fuqi. Nëse $x$ përkon me një nyje, kthehet drejtpërdrejt $f_i$ për të shmangur pjesëtimin me zero.

\subsection{Interpolimi Hermitian dhe ekstrapolimi i Richardson-it}
Interpolimi Hermitian përdor jo vetëm $f(x_i)$, por edhe derivate. Për dy nyje dhe vlera të derivateve ndërtohet një polinom kub. Në formë të përgjithshme, nyjet përsëriten në tabelën e diferencave të pjesëtuara dhe diferenca e parë në një nyje të përsëritur zëvendësohet nga $f'(x_i)$.

Ekstrapolimi nis nga një përafrim me zhvillim asimptotik
\begin{equation}
 A(h)=A+c_ph^p+c_{p+1}h^{p+1}+\cdots.
\end{equation}
Dy vlera me hapa $h$ dhe $h/r$ japin
\begin{equation}
 A_{\mathrm{R}}=\frac{r^pA(h/r)-A(h)}{r^p-1}
 =A+\Oh(h^{p+1}),
\end{equation}
ose rend edhe më të lartë kur zhvillimi përmban vetëm fuqi çifte. Kjo është ideja pas derivimit të përmirësuar dhe integrimit Romberg. Ekstrapolimi nuk duhet zbatuar jashtë regjimit asimptotik; nëse hapi është shumë i madh, modeli i gabimit nuk vlen, ndërsa për hap tepër të vogël dominon rrumbullakimi.

\subsection{Katrorët më të vegjël dhe QR}
Për modelin linear $\vect y=X\vect\beta+\vect\varepsilon$, minimizohet
\begin{equation}
 S(\vect\beta)=\|\vect y-X\vect\beta\|_2^2.
\end{equation}
Gradienti është $\nabla S=-2X^T(\vect y-X\vect\beta)$. Kushti $\nabla S=0$ jep ekuacionet normale
\begin{equation}
 X^TX\widehat{\vect\beta}=X^T\vect y.
\end{equation}
Gjeometrikisht, mbetja është ortogonale me hapësirën e kolonave të $X$. Megjithatë, $\kappa(X^TX)=\kappa(X)^2$, prandaj formimi i ekuacioneve normale përkeqëson kushtëzimin. Me $X=QR$, ku $Q^TQ=I$, zgjidhet
\begin{equation}
 R\widehat{\vect\beta}=Q^T\vect y,
\end{equation}
pa katrorëzuar numrin e kushtëzimit.

Nëse gabimet kanë kovariancë $\Sigma$, kriteri i peshuar është
\begin{equation}
 (\vect y-X\vect\beta)^T\Sigma^{-1}(\vect y-X\vect\beta).
\end{equation}
Kjo formulë tregon se peshat nuk janë dekorim: ato përfaqësojnë modelin probabilitar të matjeve. Kur gabimet kanë shpërndarje Gauss, katrorët më të vegjël përkojnë me vlerësimin e gjasës më të madhe.

\subsection{Shmangiet minimale absolute}
Kriteri
\begin{equation}
 \min_{\vect\beta}\sum_{i=1}^{n}|y_i-\vect x_i^T\vect\beta|
\end{equation}
është më pak i ndjeshëm ndaj vlerave ekstreme se kriteri kuadratik. Duke futur $u_i\ge0$ me
\begin{equation}
 -u_i\le y_i-\vect x_i^T\vect\beta\le u_i,
\end{equation}
problemi bëhet programim linear me objektiv $\min\sum_i u_i$. Funksioni absolut nuk është i derivueshëm në zero, prandaj nuk duhet trajtuar verbërisht me ekuacionet normale.

\begin{lstlisting}[language=Python,caption={Përshtatje lineare me QR dhe diagnostikë të mbetjeve.}]
import numpy as np

def pershtatje_qr(x, y):
    x = np.asarray(x, float); y = np.asarray(y, float)
    X = np.column_stack((np.ones_like(x), x))
    Q, R = np.linalg.qr(X, mode="reduced")
    beta = np.linalg.solve(R, Q.T @ y)
    mbetje = y - X @ beta
    sigma2 = (mbetje @ mbetje) / (len(y) - X.shape[1])
    kov_beta = sigma2 * np.linalg.inv(R) @ np.linalg.inv(R).T
    return beta, mbetje, kov_beta
\end{lstlisting}

\subsection{Përmbledhje dhe ushtrime}
Interpolimi supozon se të dhënat duhet të kalohen saktësisht; regresi pranon mbetje dhe kërkon një model të tyre. Zgjedhja nuk është thjesht numerike, por lidhet me kuptimin e të dhënave.
\begin{enumerate}
\item Vërtetoni formulën e mbetjes së Lagranzhit duke ndërtuar një funksion me $n+2$ zero.
\item Krahasoni nyjet e baraslarguara dhe të Çebishevit për funksionin Runge.
\item Nxirrni Richardson-in për një metodë të rendit të dytë dhe ndërtoni tabelën Romberg.
\item Përshtatni ligjin e tretë të Keplerit në shkallë lineare dhe logaritmike; diskutoni modelin e gabimit në secilin rast.
\item Formuloni regresin me shmangie minimale absolute si program linear.
\end{enumerate}

\subsection{Ndërtimi i interpoluesit Hermitian}
Për nyje $x_i$ ku njihen $f_i=f(x_i)$ dhe $f_i'=f'(x_i)$, interpoluesi Hermitian me $n+1$ nyje ka gradë jo më të madhe se $2n+1$. Duke përdorur bazën e Lagranzhit $L_i(x)$, përkufizohen
\begin{equation}
 H_i(x)=\left[1-2L_i'(x_i)(x-x_i)\right]L_i(x)^2,
 \qquad
 K_i(x)=(x-x_i)L_i(x)^2.
\end{equation}
Ato plotësojnë $H_i(x_j)=\delta_{ij}$, $H_i'(x_j)=0$, $K_i(x_j)=0$ dhe $K_i'(x_j)=\delta_{ij}$. Prandaj
\begin{equation}
 p(x)=\sum_{i=0}^{n}\left[f_iH_i(x)+f_i'K_i(x)\right].
\end{equation}
Nëse $f\in C^{2n+2}$, mbetja është
\begin{equation}
 f(x)-p(x)=\frac{f^{(2n+2)}(\xi_x)}{(2n+2)!}
 \prod_{i=0}^{n}(x-x_i)^2.
\end{equation}
Fuqia katrore tregon se vlera dhe pjerrësia përputhen në çdo nyje. Në të dhëna eksperimentale, derivatet e vlerësuara me zhurmë mund ta përkeqësojnë rezultatin; interpolimi Hermitian është më i vlefshëm kur derivatet vijnë nga ligji fizik ose maten me cilësi të mjaftueshme.

\subsubsection{Regresi, gjasa dhe pacaktueshmëria e parametrave}
Në modelin $\vect y=X\vect\beta+\vect\varepsilon$, me $\vect\varepsilon\sim\mathcal N(0,\sigma^2I)$, logaritmi i gjasës ndryshon nga një konstante me
\begin{equation}
 -\frac{1}{2\sigma^2}\|\vect y-X\vect\beta\|_2^2.
\end{equation}
Prandaj vlerësuesi i katrorëve më të vegjël është edhe vlerësues i gjasës më të madhe. Kur $X$ ka rang të plotë,
\begin{equation}
 \widehat{\vect\beta}=(X^TX)^{-1}X^T\vect y,
 \qquad
 \Cov(\widehat{\vect\beta})=\sigma^2(X^TX)^{-1}.
\end{equation}
Kjo formulë përdoret për teori, por jo domosdoshmërisht për llogaritje: QR ose SVD janë më të qëndrueshme. Varianca vlerësohet me
\begin{equation}
 s^2=\frac{\|\vect r\|^2}{n-p},
\end{equation}
ku $p$ është numri i parametrave. Pjesëtimi me $n-p$ merr parasysh shkallët e lirisë të përdorura në përshtatje.

Një interval për mesataren e modelit në $\vect x_*$ ka variancë $s^2\vect x_*^T(X^TX)^{-1}\vect x_*$. Një interval parashikimi për një matje të re shton edhe $s^2$, sepse matja e re ka zhurmën e vet. Ngatërrimi i këtyre dy intervaleve çon në deklarim tepër optimist të parashikimit.

\subsubsection{Diagnostika dhe zgjedhja e modelit}
Mbetjet duhet të kontrollohen ndaj vlerave të përshtatura, kohës dhe ndryshoreve shpjeguese. Një hark sistematik tregon jolinearitet të pambuluar; një hinkë tregon variancë jo konstante; grupe të ndara tregojnë ndryshore të munguar. Autokorrelacioni i mbetjeve shkel pavarësinë dhe ndryshon pacaktueshmërinë e parametrave.

Shtimi i parametrave ul gjithmonë shumën e katrorëve, prandaj ajo nuk mjafton për zgjedhje modeli. Vleftësimi i kryqëzuar ndan të dhënat në pjesë trajnimi dhe kontrolli. Kriteret AIC ose BIC kombinojnë përshtatjen me penalizim të kompleksitetit, por bazohen në supozime probabilitare që duhet deklaruar. Në ligjin e Keplerit, përshtatja e $T^2=a^3$ mund të bëhet me një parametër fizik; një polinom i gradës së lartë mund të përshtatë pikat më mirë, por nuk jep interpretim dhe ekstrapolon keq.

Ekstrapolimi është gjithmonë më i rrezikshëm se interpolimi sepse largohet nga mbështetja e të dhënave. Intervalet statistikë të një modeli nuk përfshijnë automatikisht gabimin nga forma e gabuar e modelit. Për këtë arsye, çdo ekstrapolim duhet të shoqërohet me arsyetim fizik dhe analizë ndjeshmërie ndaj modeleve alternative.

